\begin{figure}[H]
\centering
\begin{tikzpicture}[>=Stealth,scale=0.91]
  \draw[rounded corners,very thick,gray!65] (-5,-2.35) rectangle (5.5,2.35);
  \draw[->] (-5.4,0) -- (7.2,0) node[right] {eje óptico};
  \filldraw[fill=cyan!14,draw=blue!65,very thick]
    (-4.45,-2.05) .. controls (-3.85,-1.2) and (-3.85,1.2) .. (-4.45,2.05)
    .. controls (-4.05,1.2) and (-4.05,-1.2) .. cycle;
  \node[rotate=90] at (-4.82,0) {menisco corrector};
  \draw[very thick] (4.5,-2.0) .. controls (5.15,-1) and (5.15,-0.35) .. (4.55,-0.28);
  \draw[very thick] (4.55,0.28) .. controls (5.15,0.35) and (5.15,1) .. (4.5,2.0);
  \node[right] at (4.85,-1.55) {primario esférico};
  \draw[line width=3pt,gray!80] (-4.02,-0.5) .. controls (-3.82,-0.25) and (-3.82,0.25) .. (-4.02,0.5);
  \node[align=center] at (-1.9,1.82) {zona aluminizada\\(secundario)};
  \draw[->,thin] (-2.4,1.62) -- (-3.9,0.43);
  \foreach \y in {-1.3,1.3}{
    \draw[->,blue,thick] (-5.0,\y) -- (4.62,\y);
  }
  \draw[red,thick] (4.62,1.3) -- (-3.95,0.38);
  \draw[red,thick] (4.62,-1.3) -- (-3.95,-0.38);
  \draw[->,green!60!black,thick] (-3.95,0.38) -- (6.2,0);
  \draw[->,green!60!black,thick] (-3.95,-0.38) -- (6.2,0);
  \filldraw[fill=gray!25] (6.1,-0.5) rectangle (6.35,0.5);
  \node[above] at (6.22,0.58) {plano focal};
\end{tikzpicture}
\caption{Elementos ópticos de un telescopio Maksutov--Cassegrain.}
\end{figure}